\section{Experimental Setup}
\label{sec:setup}
\textbf{Data.} C-MAPSS FD001 (one operating condition, one fault mode; 100 training and 100 test engines) and FD002 (six operating conditions, one fault mode; 260 training and 259 test engines)~\cite{saxena2008damage}. Models are evaluated on the official test files: one window per engine (the last $W$ cycles) against the provided RUL, with the true RUL capped at 125 cycles.

\textbf{Partitioning.} Two clients. For FD001, centralised, local-only and federated runs use the same 90 training units (data-parity fix recorded in the audit log). The checked-in FD002 generator sorts the 260 engines by lifespan and assigns the shorter-lived half to client~1 and the rest to client~2, inducing label skew rather than an operating-condition split. \todo{Recover and publish the exact client unit IDs used for the archived physical FD002 results; audit notes call the runs operating-condition-partitioned, which conflicts with the checked-in generator.}

\textbf{Training.} Adam, learning rate $10^{-3}$, batch size 64, MSE loss. Federated: 2 local epochs per round, 5 rounds, FedAvg. Centralised and local-only baselines: 10 epochs. The constant baseline predicts the training-set mean RUL. \todo{Min-clients/fraction settings and server-side evaluation details from server.py}

\textbf{Seeds.} 42, 100, 2024, 7, 777 ($n=5$); client and server seeds are set from the \texttt{SEED} environment variable. Earlier FD001 federated runs fixed the client seed at 42, so their variance reflects the server seed only (audit log). \todo{State whether the reported FD001 files were produced before or after this change}

\textbf{Hardware.} Federated runs: two Raspberry Pi 4 clients plus an Ubuntu PC coordinator. \todo{Pi RAM, CPU governor, PC specs}. Centralised and local-only baselines run via \texttt{final\_benchmark.py}; \todo{which machine}.

\textbf{Statistics.} Tables report mean $\pm$ the standard deviation across seeds as produced by \texttt{final\_benchmark.py}; for the federated rows this is the population standard deviation (ddof\,=\,0). Sample standard deviation computed from the raw JSON files is \FdOneFedRmseSampleStd{} (FD001) and \FdTwoFedRmseSampleStd{} (FD002) RMSE.

\textbf{Loss function comparison protocol.} To isolate the effect of the training objective (\texttt{LOSS\_FN=mse} vs.\ \texttt{LOSS\_FN=asymmetric} with $\alpha_{\text{late}}{=}2.5$, $\alpha_{\text{early}}{=}0.5$), \texttt{benchmark\_asymmetric\_loss.py} evaluates Isolated Local (10 epochs, averaged across the two clients), FedAvg (in-process simulation, 5 rounds $\times$ 2 local epochs), and Centralised (10 epochs) over five seeds (\texttt{42, 123, 456, 789, 1024}) on both FD001 (100 test engines; odd/even 2-client unit split with healthy-regime min--max scaling from the first 50 cycles per unit) and FD002 (259 test engines; \texttt{partition.json} and \texttt{global\_scale.json}). All runs use a 90/10 unit-level train/validation split, Adam ($\text{lr}{=}10^{-3}$), batch size 64, and report Test RMSE, NASA Asymmetric Score, and late-prediction rate ($\hat{y} > y$) to \texttt{results/asymmetric\_loss\_comparison.json}.
